% TOP_PROBLEM: {"id": 2065, "title": "Erdős Problem #346", "category_id": 1, "category": {"id": 1, "name": "number_theory", "display_name": "Number Theory", "description": "Properties of integers, prime numbers, Diophantine equations.", "slug": "number-theory", "order_index": 1, "created_at": "2026-07-31T15:26:25.670Z"}, "status": "open"}
% FIRST_POSTED: null
% ENTRY_KIND: statement_only
% Imported from: unsolvedmath_additional_500.tex; UnsolvedMath ID 2065
% !TeX program = lualatex
% Compile twice: lualatex 2065.tex
% Self-contained: no external bibliography, images, or \input files.
% Numeric section labels and visible numbers are original UnsolvedMath IDs.
\documentclass[11pt,a4paper]{article}
\usepackage[margin=24mm,headheight=14pt]{geometry}
\usepackage{fontspec}
\setmainfont{Latin Modern Roman}
\setsansfont{Latin Modern Sans}
\setmonofont{Latin Modern Mono}
\usepackage{amsmath,amssymb,mathtools}
\usepackage{unicode-math}
\setmathfont{Latin Modern Math}
\usepackage{microtype}
\usepackage{enumitem,xurl,needspace}
\usepackage[dvipsnames]{xcolor}
\usepackage{hyperref}
\hypersetup{unicode=true,colorlinks=true,linkcolor=MidnightBlue,urlcolor=MidnightBlue,
 pdftitle={UnsolvedMath Problem 2065},pdfauthor={Source-annotated compilation},bookmarksnumbered=true}
\usepackage{bookmark,tocloft,titlesec,fancyhdr}
\setlength{\cftsecnumwidth}{4.2em}
\cftsetpnumwidth{2.5em}
\cftsetrmarg{4em}
\titleformat{\section}{\Large\bfseries\raggedright}{\thesection}{0.7em}{}
\pagestyle{fancy}\fancyhf{}
\fancyhead[L]{\small\sffamily Additional UnsolvedMath statements}
\fancyhead[R]{\small Original catalogue IDs}
\fancyfoot[C]{\thepage}
\setlength{\parindent}{0pt}
\setlength{\parskip}{5pt plus 1pt}
\setlength{\emergencystretch}{3em}
\setcounter{tocdepth}{1}\setcounter{secnumdepth}{1}
\setlist[itemize]{leftmargin=3em,itemsep=4pt,topsep=4pt,parsep=0pt}
\allowdisplaybreaks
\newcommand{\N}{\mathbb{N}}
\newcommand{\Z}{\mathbb{Z}}
\newcommand{\Q}{\mathbb{Q}}
\newcommand{\R}{\mathbb{R}}
\newcommand{\C}{\mathbb{C}}
\newcommand{\F}{\mathbb{F}}
\newcommand{\A}{\mathbb{A}}
\newcommand{\PP}{\mathbb{P}}
\newcommand{\E}{\mathbb{E}}
\newcommand{\eps}{\varepsilon}
\newcommand{\dd}{\,\mathrm d}
\newcommand{\sref}[2]{\hyperlink{source:#1:#2}{[S#2]}}
\newif\ifshowcatalogue
\showcataloguetrue
\newcommand{\cataloguescope}[1]{\ifshowcatalogue
 \paragraph{Statement recorded in the supplied source.}
 #1\fi}
\begin{document}
% UnsolvedMath ID 2065; source catalogue code EP-346

\setcounter{section}{2064}

\section{A Golden-Ratio Law for Critically Complete Sequences}\label{2065}

{\small\sffamily UnsolvedMath ID 2065 · Number Theory}

\subsection{Definitions and mathematical statement}

\paragraph{Shared notation.}Unless a section says otherwise, $\N=\{1,2,\ldots\}$,
$\Z$ is the ring of integers, $\Q,\R,\C$ are the rational, real, and complex
fields, $[N]=\{1,\ldots,\lfloor N\rfloor\}$, and $\log$ is the natural
logarithm. For real functions of a parameter tending to infinity,
$f\ll g$ and $f=O(g)$ mean $|f|\leq Cg$ eventually for some fixed $C$;
$f\gg g$ reverses this comparison; $f=o(g)$ means $f/g\to0$;
$f\sim g$ means $f/g\to1$; and $f\asymp g$ means both order bounds.
A subscript on $O$ or $\ll$ specifies allowed constant dependence.
The expression $n^{o(1)}$ in an upper bound means growth smaller than
$n^\epsilon$ for each fixed $\epsilon>0$. Local definitions govern any
exception, finite-field convention, or use of the same symbol for another
object. An asymptotic claim is not automatically an effective algorithm.



Completeness is unchanged by ignoring finitely many exceptional target integers. The first deletion property quantifies over every finite subset of $A$, while the second quantifies over every infinite subset. The lower ratio bound uses one fixed positive $\epsilon$. \sref{2065}{1} \sref{2065}{2}

\paragraph{Statement recorded in the supplied source.}
Let $A=\{1\leq a_1< a_2<\cdots\}$ be a set of integers such that
\begin{itemize}
\item  $A\backslash B$ is complete for any finite subset $B$ and 
\item  $A\backslash B$ is not complete for any infinite subset $B$.
\end{itemize}
(Here 'complete' means all sufficiently large integers can be written as a sum of distinct members of the sequence.)
Is it true that if $a_{n+1}/a_n \geq 1+\epsilon$ for some $\epsilon>0$ and all $n$ then $ \lim_n \frac{a_{n+1}}{a_n}=\frac{1+\sqrt{5}}{2}? $ \sref{2065}{1}

{\small\emph{Transcription note.} Only unambiguous escape/formatting damage and nonmathematical serialization debris were repaired in the displayed source statement.}

\subsection{Short English statement}

If completeness survives every finite deletion but no infinite deletion, must uniformly separated successive ratios converge to the golden ratio?

\subsection{Sources}

{\small

\begin{itemize}

\item[\hypertarget{source:2065:1}{S1}] ULAM.AI, \emph{UnsolvedMath}, supplied \texttt{problems.json}, record 2065 (EP-346), statement and background. The embedded literature-review date is 2026-08-17; that is the archive’s review date, not a fresh certification. Dataset landing page: \url{https://huggingface.co/datasets/ulamai/UnsolvedMath}.

\item[\hypertarget{source:2065:2}{S2}] T. F. Bloom, \emph{Erdős Problems}, Problem 346, \url{https://www.erdosproblems.com/346}. Reference imported from the supplied record; the live problem page was not independently checked for this compilation.

\item[\hypertarget{source:2065:3}{S3}] Erd\H{o}s, P. and Graham, R., Old and new problems and results in combinatorial number theory. Monographies de L'Enseignement Mathematique (1980). Bibliographic information imported from the supplied record.

\item[\hypertarget{source:2065:4}{S4}] Graham, R. L., A property of Fibonacci numbers. Fibonacci Quart. (1964), 1-10. Bibliographic information imported from the supplied record.

\end{itemize}

}

\label{end:2065}
\end{document}
